\documentclass[11pt,a4paper]{article}

\usepackage[margin=1in]{geometry}
\usepackage{amsmath,amssymb,mathtools}
\usepackage[most]{tcolorbox}
\usepackage{enumitem}
\usepackage{hyperref}
\usepackage{fancyhdr}
\usepackage{braket}
\usepackage[]{cancel}
\usepackage[]{tikz}

\setlength{\parindent}{0pt}
\setlength{\parskip}{0.6em}

% ---------- Page style ----------
\pagestyle{fancy}
\fancyhf{}
\lhead{Advanced Statistical Mechanics / Statistical Mechanics I}
\rhead{Solution Manual II}
\cfoot{\thepage}

% ---------- Hyperref ----------
\hypersetup{
	colorlinks=true,
	linkcolor=blue,
	urlcolor=blue,
	citecolor=blue
}

% ---------- Counters ----------
\newcounter{problem}
\newcounter{saveequation}

% Rounded box for each problem, with automatic numbering
\newtcolorbox{problembox}{
	enhanced,
	breakable,
	colback=white,
	colframe=black,
	boxrule=0.8pt,
	arc=2.5mm,
	left=2mm,
	right=2mm,
	top=1.5mm,
	bottom=1.5mm,
	before upper={\refstepcounter{problem}\textbf{(\theproblem)}\quad\ignorespaces}
}

% Plain solution environment
% Equations inside are numbered as problem.equation, e.g. 3.1, 3.2
% Equation numbering outside remains unchanged
\newenvironment{solution}
{
	\par
	\setcounter{saveequation}{\value{equation}}
	\begingroup
	\setcounter{equation}{0}
	\renewcommand{\theequation}{\arabic{problem}.\arabic{equation}}
	\renewcommand{\theHequation}{sol.\arabic{problem}.\arabic{equation}}
	\noindent\textbf{Solution.}\quad
}
{
	\par
	\endgroup
	\setcounter{equation}{\value{saveequation}}
}

% Background matter box
\definecolor{bgmatterback}{RGB}{245,248,252}
\definecolor{bgmatterframe}{RGB}{120,140,170}
\definecolor{bgmattertitle}{RGB}{70,92,125}

\newtcolorbox{backgroundmatter}[1][]{
	enhanced,
	breakable,
	colback=bgmatterback,
	colframe=bgmatterframe,
	coltitle=white,
	colbacktitle=bgmattertitle,
	boxrule=0.9pt,
	arc=2.8mm,
	left=2.5mm,
	right=2.5mm,
	top=1.5mm,
	bottom=1.5mm,
	title=Background Matter,
	fonttitle=\bfseries,
	attach boxed title to top left={xshift=2mm,yshift=-2mm},
	boxed title style={
		arc=2mm,
		boxrule=0pt,
		left=1.5mm,
		right=1.5mm,
		top=0.8mm,
		bottom=0.8mm
	},
	before upper={
		\setlength{\parindent}{2em}
		\setlength{\parskip}{0pt}
	},
	#1
}

% Final answer box
\definecolor{finalansback}{RGB}{248,251,247}
\definecolor{finalansframe}{RGB}{88,122,96}
\definecolor{finalanstitle}{RGB}{63,96,70}

\newtcolorbox{finalanswer}[1][]{
	enhanced,
	breakable,
	colback=finalansback,
	colframe=finalansframe,
	coltitle=white,
	colbacktitle=finalanstitle,
	boxrule=1pt,
	arc=2.8mm,
	left=2.5mm,
	right=2.5mm,
	top=1.5mm,
	bottom=1.5mm,
	title=Final Answer,
	fonttitle=\bfseries,
	attach boxed title to top left={xshift=2mm,yshift=-2mm},
	boxed title style={
		arc=2mm,
		boxrule=0pt,
		left=1.5mm,
		right=1.5mm,
		top=0.8mm,
		bottom=0.8mm
	},
	#1
}

\title{\textbf{Report II}\\[0.4em]
	\large Advanced Statistical Mechanics / Statistical Mechanics I}
\author{Bufan Zheng\\[0.2em]\normalsize Student ID: 35-256419}
\date{May 25, 2026}

\begin{document}
	
	\maketitle
	
	The final results can be found in \eqref{fin:1}, \eqref{fin:2}, \eqref{fin:3}, \eqref{fin:4}, \eqref{fin:5}, \eqref{fin:6}, \eqref{fin:7} and \eqref{fin:8}.
	
	\section{Gaussian integral}
	
	Let $\boldsymbol{x} = (x_1,\cdots,x_N)$ be an $N$-dimensional real vector, and let $A$ be an $N\times N$ real, symmetric, and positive-definite matrix.
	
	\begin{problembox}
		Calculate the Gaussian integral
		\begin{equation}
			\int_{-\infty}^{\infty} d^N x\, e^{-\frac{1}{2}\boldsymbol{x}^{\mathrm{T}}A\boldsymbol{x}}
			\left(=
			\int_{-\infty}^{\infty} dx_1\cdots\int_{-\infty}^{\infty} dx_N\,
			e^{-\frac{1}{2}\sum_{i,j=1}^{N}x_iA_{ij}x_j}
			\right).
		\end{equation}
	\end{problembox}
	\begin{solution}
		Since $A$ is a real symmetric $N\times N$ matrix, we can diagonalize it using an orthogonal matrix,
		\begin{equation}
			P^\mathrm{T} A P = \operatorname{diag}(a_1,a_2,\cdots,a_N),\quad P\in SO(N)
		\end{equation}
		Using the substitution $\mathbf{x}\to P\cdot\mathbf{x}$, we obtain
		\begin{equation}
				\int_{-\infty}^{\infty} d^N x\, e^{-\frac{1}{2}\boldsymbol{x}^{\mathrm{T}}A\boldsymbol{x}}
				=
				\frac{1}{\det P}
				\int_{-\infty}^{\infty} dx_1\cdots\int_{-\infty}^{\infty} dx_N\,
				e^{-\frac{1}{2}\sum_{i=1}^{N}a_i x_i^2}=\prod_{i=1}^{N}\int_{-\infty}^{\infty} dx_i e^{-\frac12 a_i x_i^2}
		\end{equation}
		Using the following formula,
		\begin{equation}
			\int_{-\infty}^{\infty}dx e^{-ax^2}=\sqrt{\frac\pi a}
		\end{equation}
		We obtain the final answer:
		\begin{finalanswer}
			\begin{equation}
				\label{fin:1}
				\prod_{i=1}^{N}\sqrt{\frac{2\pi}{a_i}}=\sqrt{\frac{(2\pi)^{N}}{\det A}}
			\end{equation}
		\end{finalanswer}
	\end{solution}
	
	\begin{problembox}
		For an $N$-dimensional real vector $\boldsymbol{b}=(b_1,\cdots,b_N)$, calculate the Gaussian integral
		\begin{equation}
			\int_{-\infty}^{\infty} d^N x\, e^{-\frac{1}{2}\boldsymbol{x}^{\mathrm{T}}A\boldsymbol{x}+\boldsymbol{b}^{\mathrm{T}}\boldsymbol{x}}.
		\end{equation}
	\end{problembox}
	\begin{solution}
		Diagonalizing $A$ as in the previous problem and denoting $P^\mathrm{T}\cdot\boldsymbol{b}$ by $(b'_1,\cdots,b'_N)$, we have
		\begin{equation}
			\begin{aligned}
				\int_{-\infty}^{\infty} d^N x\, e^{-\frac{1}{2}\boldsymbol{x}^{\mathrm{T}}A\boldsymbol{x}+\boldsymbol{b}^{\mathrm{T}}\boldsymbol{x}}&=\int_{-\infty}^{\infty} dx_1\cdots\int_{-\infty}^{\infty} dx_N\,
				e^{-\frac{1}{2}\sum_{i=1}^{N}a_i x_i^2+b'_i x_i}\\
				&=\prod_{i=1}^{N}\int dx_i e^{-\frac12 a_i x_i^2+b'_i x_i}=\prod_{i=1}^{N}\int dx_i e^{-\frac12 a_i (x_i-\frac{b'_i}{a_i})^2}\cdot e^{\frac{{b'_i}^2}{2a_i}}\\
				&=\sqrt{\frac{(2\pi)^{N}}{\det A}}\cdot e^{\sum_{i=1}^N\frac{{b'_i}^2}{2a_i}}
			\end{aligned}
		\end{equation}
		Notice that
		\begin{equation}
			\sum_{i=1}^N\frac{{b'_i}^2}{2a_i}=\frac12 (P^\mathrm{T}\cdot\boldsymbol{b})^\mathrm{T}\cdot (P^\mathrm{T} A P)^{-1}\cdot P^\mathrm{T}\cdot\boldsymbol{b}=\frac12 \boldsymbol{b}^\mathrm{T} A^{-1} \boldsymbol{b}
		\end{equation}
		Then the final answer is:
		\begin{finalanswer}
			\begin{equation}
				\label{fin:2}
				\sqrt{\frac{(2\pi)^{N}}{\det A}}\cdot e^{\frac12 \boldsymbol{b}^\mathrm{T} A^{-1} \boldsymbol{b}}
			\end{equation}
		\end{finalanswer}
	\end{solution}
	
	For a function $f(\boldsymbol{x})$ of the $N$-dimensional vector $\boldsymbol{x}$, we introduce the average over the Gaussian distribution:
	\begin{equation}
		\left\langle f(\boldsymbol{x}) \right\rangle :=
		\frac{\displaystyle \int_{-\infty}^{\infty} d^N x\, f(\boldsymbol{x})\, e^{-\frac{1}{2}\boldsymbol{x}^{\mathrm{T}}A\boldsymbol{x}}}
		{\displaystyle \int_{-\infty}^{\infty} d^N x\, e^{-\frac{1}{2}\boldsymbol{x}^{\mathrm{T}}A\boldsymbol{x}}}.
	\end{equation}
	
	\begin{problembox}
		Calculate $\left\langle x_i x_j \right\rangle$ $(i,j\in\{1,\cdots,N\})$.
	\end{problembox}
	\begin{solution}
		With the following notation:
		\begin{equation}
			Z[\boldsymbol{b}]:=\int_{-\infty}^{\infty} d^N x\, e^{-\frac{1}{2}\boldsymbol{x}^{\mathrm{T}}A\boldsymbol{x}+\boldsymbol{b}^{\mathrm{T}}\boldsymbol{x}}
		\end{equation}
		It is easy to show that
		\begin{equation}
			\begin{aligned}
				& \braket{x_i}:=\frac{\displaystyle \int_{-\infty}^{\infty} d^N x\, x_i\, e^{-\frac{1}{2}\boldsymbol{x}^{\mathrm{T}}A\boldsymbol{x}}}
				{\displaystyle \int_{-\infty}^{\infty} d^N x\, e^{-\frac{1}{2}\boldsymbol{x}^{\mathrm{T}}A\boldsymbol{x}}}=\frac{1}{Z[0]}\left.\frac{\partial}{\partial b_i}Z[\boldsymbol{b}]\right|_{\boldsymbol{b}=0}\\
				& \braket{x_ix_j}:=\frac{\displaystyle \int_{-\infty}^{\infty} d^N x\, x_ix_j\, e^{-\frac{1}{2}\boldsymbol{x}^{\mathrm{T}}A\boldsymbol{x}}}
				{\displaystyle \int_{-\infty}^{\infty} d^N x\, 	e^{-\frac{1}{2}\boldsymbol{x}^{\mathrm{T}}A\boldsymbol{x}}}=\frac{1}{Z[0]}\left.\frac{\partial}{\partial b_i}\frac{\partial}{\partial b_j}Z[\boldsymbol{b}]\right|_{\boldsymbol{b}=0}\\
				&\cdots\\
				& \braket{x_{i_1}\cdots x_{i_n}}:=\frac{\displaystyle \int_{-\infty}^{\infty} d^N x\, x_{i_1}\cdots x_{i_n}\, e^{-\frac{1}{2}\boldsymbol{x}^{\mathrm{T}}A\boldsymbol{x}}}
				{\displaystyle \int_{-\infty}^{\infty} d^N x\, 	e^{-\frac{1}{2}\boldsymbol{x}^{\mathrm{T}}A\boldsymbol{x}}}=\frac{1}{Z[0]}\left.\frac{\partial}{\partial b_{i_1}}\cdots\frac{\partial}{\partial b_{i_n}}Z[\boldsymbol{b}]\right|_{\boldsymbol{b}=0}\\
			\end{aligned}
		\end{equation}
	Focusing on $\braket{x_ix_j}$ and using \eqref{fin:2}, we obtain:
	\begin{equation}
		\begin{aligned}
			\braket{x_ix_j}&=\frac{1}{Z[0]}\left.\frac{\partial}{\partial b_i}\frac{\partial}{\partial b_j}Z[\boldsymbol{b}]\right|_{\boldsymbol{b}=0}=\cancelto{1}{\sqrt{\frac{(2\pi)^{N}}{\det A}}\frac{1}{Z[0]}}\left.\frac{\partial}{\partial b_i}\frac{\partial}{\partial b_j} e^{\frac12 \boldsymbol{b}^\mathrm{T} A^{-1} \boldsymbol{b}}\right|_{\boldsymbol{b}=0}\\
			& =\left.e^{\frac12 \boldsymbol{b}^\mathrm{T} A^{-1} \boldsymbol{b}}\frac{\partial}{\partial b_i}\frac{\partial}{\partial b_j} \left({\frac12 \boldsymbol{b}^\mathrm{T} A^{-1} \boldsymbol{b}}\right)\right|_{\boldsymbol{b}=0}+
			\cancelto{0}{\left[\left.\frac{\partial}{\partial b_i}e^{\frac12 \boldsymbol{b}^\mathrm{T} A^{-1} \boldsymbol{b}}\frac{\partial}{\partial b_j} \left({\frac12 \boldsymbol{b}^\mathrm{T} A^{-1} \boldsymbol{b}}\right)\right|_{\boldsymbol{b}=0}+(i\leftrightarrow j)\right]}\\
			&= \frac{1}{2}\left(\left(A^{-1}\right)_{ij}+\left(A^{-1}\right)_{ji}\right)
		\end{aligned}
	\end{equation}
	Since $A$ is symmetric, so is $A^{-1}$. Therefore, the final answer is:
	\begin{finalanswer}
		\begin{equation}
		\label{fin:3}
		\braket{x_ix_j}=\left(A^{-1}\right)_{ij}
		\end{equation}
	\end{finalanswer}
	\end{solution}
	\begin{problembox}
		Show
		\begin{equation}
			\left\langle x_i x_j x_k x_l \right\rangle
			= \left\langle x_i x_j \right\rangle \left\langle x_k x_l \right\rangle
			+ \left\langle x_i x_k \right\rangle \left\langle x_j x_l \right\rangle
			+ \left\langle x_i x_l \right\rangle \left\langle x_j x_k \right\rangle
			\quad (i,j,k,l\in\{1,\cdots,N\}).
		\end{equation}
	\end{problembox}
	\begin{solution}
		The proof is straightforward by definition. Notice that the only contribution comes from the second-order term in the Taylor expansion of 
		$e^{\frac12 \boldsymbol{b}^\mathrm{T} A^{-1} \boldsymbol{b}}$, which is the fourth-order term in $\boldsymbol{b}$.
	\begin{finalanswer}
				\begin{equation}
				\label{fin:4}
				\begin{aligned}
					\braket{x_ix_jx_kx_l}
					&=\frac{1}{Z[0]}
					\left.
					\frac{\partial}{\partial b_i}
					\frac{\partial}{\partial b_j}
					\frac{\partial}{\partial b_k}
					\frac{\partial}{\partial b_l}
					Z[\boldsymbol{b}]
					\right|_{\boldsymbol{b}=0}\\
					&=\cancelto{1}{\sqrt{\frac{(2\pi)^N}{\det A}}\frac{1}{Z[0]}}
					\left.
					\frac{\partial}{\partial b_i}
					\frac{\partial}{\partial b_j}
					\frac{\partial}{\partial b_k}
					\frac{\partial}{\partial b_l}
					e^{\frac12 \boldsymbol{b}^{\mathrm T}A^{-1}\boldsymbol{b}}
					\right|_{\boldsymbol{b}=0}\\
					&=
					\left.
					\frac{\partial}{\partial b_i}
					\frac{\partial}{\partial b_j}
					\frac{\partial}{\partial b_k}
					\frac{\partial}{\partial b_l}
					\left[
					\frac{1}{2!}
					\left(
					\frac12 \boldsymbol{b}^{\mathrm T}A^{-1}\boldsymbol{b}
					\right)^2
					\right]
					\right|_{\boldsymbol{b}=0}\\
					&=
					\left.
					\frac{\partial}{\partial b_i}
					\frac{\partial}{\partial b_j}
					\frac{\partial}{\partial b_k}
					\frac{\partial}{\partial b_l}
					\left[
					\frac18
					\sum_{m,n,p,q=1}^{N}
					\left(A^{-1}\right)_{mn}
					\left(A^{-1}\right)_{pq}
					b_m b_n b_p b_q
					\right]
					\right|_{\boldsymbol{b}=0}\\
					&=
					\frac18
					\sum_{m,n,p,q=1}^{N}
					\left(A^{-1}\right)_{mn}
					\left(A^{-1}\right)_{pq}
					\left.
					\frac{\partial}{\partial b_i}
					\frac{\partial}{\partial b_j}
					\frac{\partial}{\partial b_k}
					\frac{\partial}{\partial b_l}
					b_m b_n b_p b_q
					\right|_{\boldsymbol{b}=0}\\
					&=
					\frac18\cdot 8
					\left[
					\left(A^{-1}\right)_{ij}\left(A^{-1}\right)_{kl}
					+
					\left(A^{-1}\right)_{ik}\left(A^{-1}\right)_{jl}
					+
					\left(A^{-1}\right)_{il}\left(A^{-1}\right)_{jk}
					\right]\\
					&=
					\left(A^{-1}\right)_{ij}\left(A^{-1}\right)_{kl}
					+
					\left(A^{-1}\right)_{ik}\left(A^{-1}\right)_{jl}
					+
					\left(A^{-1}\right)_{il}\left(A^{-1}\right)_{jk}\\
					&=
					\braket{x_ix_j}\braket{x_kx_l}
					+
					\braket{x_ix_k}\braket{x_jx_l}
					+
					\braket{x_ix_l}\braket{x_jx_k}.
				\end{aligned}
			\end{equation}
	\end{finalanswer}
	\end{solution}
	
	\section{Gaussian theory}
	
	Let us consider the Gaussian theory in $d < 4$ dimensions:
	\begin{equation}
		Z = \int \mathcal{D}\phi\, e^{-H[\phi]},
	\end{equation}
	\begin{equation}
		H[\phi] = \int d^d r\left[\frac{1}{2}(\nabla\phi)^2 + t\phi^2\right]\qquad (t\in\mathbb{R}).
	\end{equation}
	
	\begin{problembox}
		Perform the Gaussian integral explicitly, and then write down the free energy density.
		
		[Note] You may focus solely on the most singular contribution to the free energy density and omit the $t$-independent constant term.
	\end{problembox}
	\begin{solution}
		Integrating by parts, we have:
		\begin{equation}
			\label{eq:5.1}
			Z = \int \mathcal{D}\phi\, e^{-H[\phi]}=\int \mathcal{D}\phi\, \exp\left[-\frac12\int \phi\left(-\nabla^2+2t\right)\phi\right]\propto \left[\det \left(-\nabla^2+2t\right)\right]^{-\frac12}
		\end{equation}
		By definition, and omitting the constant term, we obtain
		\begin{equation}
			F=-\frac{1}{\beta}\log Z=\frac{T}{2}\log\det  \left(-\nabla^2+2t\right)
		\end{equation}
		Since $e^{i\mathbf{k}\cdot\mathbf{r}}$ is an eigenvector of $\left(-\nabla^2+2t\right)$:
		\begin{equation}
			\left(-\nabla^2+2t\right)e^{i\mathbf{k}\cdot\mathbf{r}}=\left(\mathbf{k}^2+2t\right)e^{i\mathbf{k}\cdot\mathbf{r}}
		\end{equation}
		Therefore, 
		\begin{equation}
			\det \left(-\nabla^2+2t\right)=\prod_{\mathbf{k}\in\mathbb{R}^d}\left(\mathbf{k}^2+2t\right)
		\end{equation}
		The free energy density can be expressed as:
		\begin{equation}
			f:=\frac FV= \frac{T}{2V}\sum_{\mathbf{k}\in\mathbb{R}^d}\log\left(\mathbf{k}^2+2t\right)
			\overset{V\to\infty}{\rightsquigarrow}\frac{T}{2\cancel{V}}\cdot\frac{\cancel{V}}{(2\pi)^d}\int d^d\mathbf{k} \log\left(\mathbf{k}^2+2t\right)
		\end{equation}
		To evaluate this integral, we first differentiate it:
		\begin{equation}
			\frac{\partial f}{\partial(2t)}
			=
			\frac T2\int \frac{d^d\mathbf{k}}{(2\pi)^d}\frac{1}{k^2+2t}
			=
			\frac T2\frac{1}{(4\pi)^{d/2}}
			\Gamma\left(1-\frac d2\right)(2t)^{d/2-1}.
		\end{equation}
		Then, integrating both sides and dropping the constant term, we obtain the final result:
		\begin{finalanswer}
			\begin{equation}
				\label{fin:5}
				f=-\frac{T}{2}\frac{1}{(4\pi)^{d/2}}\Gamma\left(-\frac d2\right)(2t)^{d/2}
			\end{equation}
		\end{finalanswer}
	For odd dimensions, there is no singularity, but in even dimensions, such as $d=2$, we need to use dimensional regularization. Using the standard expansion:
	\begin{equation}
		\Gamma\left(-1+\frac\epsilon2\right)=-\left(\frac2\epsilon+1-\gamma+\mathcal{O}(\epsilon)\right),\quad \gamma =0.577216\cdots
	\end{equation}
	Setting $d=2-\epsilon$, we obtain the regularized free energy density in two dimensions:
	\begin{equation}
		f^\text{ren}_{d=2}=\frac T2\mu^\epsilon\left(\frac{t}{2\pi}\right)^{1-\epsilon/2}\left(\frac2\epsilon+1-\gamma+\mathcal{O}(\epsilon)\right)=
		\frac{Tt}{4\pi}
		\left[
		\frac{2}{\epsilon}
		+1-\gamma
		+\log(4\pi)
		+2\log\mu
		-\log(2t)
		\right]
		+O(\epsilon)
	\end{equation}
	Here $\mu^\epsilon$ is the renormalization scale introduced to compensate for the dimension of $f$ in $2-\epsilon$ dimensions.
	\end{solution}
	\begin{problembox}
		For $t > 0$, calculate the most singular contribution to the specific heat $c_+$ and obtain the critical exponent $\alpha_+$:
		\begin{equation}
			c_+ = a_+ t^{-\alpha_+}.
		\end{equation}
	\end{problembox}
	\begin{solution}
		We focus on the most singular term. Notice that $t=T-T_c$, therefore
		\begin{equation}
			\label{eq:6.1}
			\begin{aligned}
				c_+&=-T\frac{\partial^2}{\partial T^2}f\sim \frac{T_c^2}{2}\frac{1}{(4\pi)^{d/2}}\Gamma\left(-\frac d2\right) \frac{\partial}{\partial t^2} (2t)^{d/2}\\
				&=\frac{T_c^2}{2}\frac{\Gamma{\left(2-\frac d2\right)}}{(2\pi)^{d/2}}t^{\frac d2-2}
			\end{aligned}
		\end{equation}
		The critical exponent can be read off from the equation above:
		\begin{finalanswer}
			\begin{equation}
				\label{fin:6}
				a_+=\frac{T_c^2}{2}\frac{\Gamma{\left(2-\frac d2\right)}}{(2\pi)^{d/2}},\quad \alpha_+=\frac{4-d}{2}
			\end{equation}
		\end{finalanswer}
	\end{solution}
	
	\begin{problembox}
		For $t < 0$, the above Gaussian path integral becomes ill defined. To regularize the theory, let us introduce a stabilizing quartic interaction and analyze Gaussian fluctuations around a stationary configuration:
		\begin{equation}
			\label{eq:8}
			H[\phi] = \int d^d r\left[\frac{1}{2}(\nabla\phi)^2 + t\phi^2 + u\phi^4\right]\qquad (u > 0).
		\end{equation}
		On the basis of the mean-field approximation, determine the stationary point $\phi = \phi_c$ and expand the Hamiltonian $H[\phi]$ around it to the quadratic order [i.e., up to $O((\phi-\phi_c)^2)$].
	\end{problembox}
	\begin{solution}
	We decompose the quantum field into a mean-field part $\phi_c$ (with $\nabla\phi=0$) and a quantum fluctuation part $\eta$:
	\begin{equation}
		\phi(\mathbf{r})=\phi_c+\eta(\mathbf{r})
	\end{equation}
	From \eqref{eq:8}, the stationary point satisfies the equation:
	\begin{equation}
		-\nabla^2\phi+2t\phi+4u\phi^3=0
	\end{equation}
	For the mean field, this simplifies to:
	\begin{equation}
		2t\phi_c+4u\phi_c^3=0\Rightarrow\phi_c=0,\quad \phi_c^2=-\frac{t}{2u}
	\end{equation}
	It is easy to see that, for $t<0$, only the two solutions corresponding to the latter equation are stable stationary points, namely
	\begin{equation}
		\phi_c=\pm\sqrt{-\frac{t}{2u}}
	\end{equation}
	Substituting this into \eqref{eq:8} and expanding in $\eta$, we obtain:
\begin{finalanswer}
		\begin{equation}
		\label{fin:7}
		\begin{aligned}
			H[\phi_c+\eta] = &\int d^d r\left[\frac{1}{2}(\cancelto{0}{\nabla\phi_c}+\nabla\eta)^2 + t(\phi_c+\eta)^2 + u(\phi_c+\eta)^4\right]\\
			= &\int d^d r\left[\frac{1}{2}(\nabla\eta)^2 +\cancelto{ 0}{\left(2t\phi_c+4u\phi_c^3\right)}\eta+\frac12 \left(2t+12u\phi_c^2\right)\eta^2+\mathcal{O}(\eta^3)\right]\\
			&+\int d^d{r~}\left(t\phi_c^2+u\phi_c^4\right)\\
			=&\frac12 \int d^dr\eta\left(-\nabla^2-4t\right)\eta+\mathcal{O}\left((\phi-\phi_c)^3\right)+H_0
		\end{aligned}
	\end{equation}
\end{finalanswer}
	In the last step, we have used integration by parts.
	\end{solution}
	
	\begin{problembox}
		For $t < 0$, calculate the most singular contribution to the specific heat $c_-$ and obtain the critical exponent $\alpha_-$:
		\begin{equation}
			c_- = a_- (-t)^{-\alpha_-}.
		\end{equation}
		Additionally, calculate the ratio $a_+/a_-$.
		
		[Note] While the coefficients $a_\pm$ themselves are not universal, their ratio $a_+/a_-$ is universal.
	\end{problembox}
	\begin{solution}
	Comparing \eqref{eq:5.1} with \eqref{fin:7}, we find that we only need to replace $t$ in \eqref{fin:6} by $-2t$ to obtain the free energy density under the new Hamiltonian:
	\begin{equation}
		f_-=-\frac T2\frac{1}{(4\pi)^{d/2}}\Gamma\left(-\frac d2\right)(-4t)^{d/2}
	\end{equation}
	Following the calculation in \eqref{eq:6.1}, we obtain:
	\begin{equation}
		\begin{aligned}
			c_-&=-T\frac{\partial^2}{\partial T^2}f_-\sim \frac{T_c^2}{2}\frac{1}{(4\pi)^{d/2}}\Gamma\left(-\frac d2\right) \frac{\partial}{\partial (-t)^2} (-4t)^{d/2}\\
			&=\frac{T_c^2}{2}\frac{\Gamma{\left(2-\frac d2\right)}}{(2\pi)^{d/2}}\cdot 2^{d/2}\cdot t^{\frac d2-2}
		\end{aligned}
	\end{equation}
	Thus, we obtain the critical exponent:
	\begin{finalanswer}
		\begin{equation}
			\label{fin:8}
			\alpha_-=\frac{4-d}{2},\quad a_-\frac{T_c^2}{2}\frac{\Gamma{\left(2-\frac d2\right)}}{(2\pi)^{d/2}}\cdot 2^{d/2},\quad\frac{a_+}{a_-}=2^{-d/2}
		\end{equation}
	\end{finalanswer}
\end{solution}
\end{document}
